% !PRINTED 84
\weberpage{84}
%<<CONT
then $q$ is a quantity whose absolute value is a proper fraction. The \gesperrt{convergent}
infinite product
\[
P(v) = \prod_{1,\infty}^{\nu} (1 + q^{2\nu-1}e^{2\pi i v})(1 + q^{2\nu-1}e^{-2\pi i v})
\]
therefore vanishes at \gesperrt{all and only} those points at which $\vartheta_{00}(v)$
vanishes, and moreover
\[
\begin{gathered}
P(v+1) = P(v)\\
P(v+\omega) = \frac{1+q^{-1}e^{-2\pi i v}}{1+qe^{2\pi i v}}P(v) = q^{-1}e^{-2\pi i v}P(v).
\end{gathered}
\]
But by § 20, (1) this is the period property of the function $\vartheta_{00}(v)$, and
we therefore have:
\weq{1}{\vartheta_{00}(v) = Q\prod_{1,\infty}^{\nu}(1 + q^{2\nu-1}e^{2\pi i v})(1 + q^{2\nu-1}e^{-2\pi i v}),}
where $Q$ is a factor independent of $v$. By the formulae (8) of § 21 one obtains from
this, on putting in place of $v$
\[
v + \frac{1}{2}, \qquad v + \frac{\omega}{2}, \qquad v + \frac{1+\omega}{2}
\]
respectively:
\weq{2}{\vartheta_{01}(v) = Q\prod_{1,\infty}^{\nu}(1 - q^{2\nu-1}e^{2\pi i v})(1 - q^{2\nu-1}e^{-2\pi i v}),}
\weq{3}{\vartheta_{10}(v) = Qq^{\frac{1}{4}}(e^{\pi i v} + e^{-\pi i v})\prod_{1,\infty}^{\nu}(1 + q^{2\nu}e^{2\pi i v})(1 + q^{2\nu}e^{-2\pi i v}),}
\weq{4}{\vartheta_{11}(v) = -iQq^{\frac{1}{4}}(e^{\pi i v} - e^{-\pi i v})\prod_{1,\infty}^{\nu}(1 - q^{2\nu}e^{2\pi i v})(1 - q^{2\nu}e^{-2\pi i v}).}
If one sets $v = 0$ in these formulae, in the last one after a single differentiation,
there follows
\weq{5}{\begin{gathered}
\vartheta_{00} = Q\prod_{1,\infty}^{\nu}(1 + q^{2\nu-1})^2,\\
\vartheta_{01} = Q\prod_{1,\infty}^{\nu}(1 - q^{2\nu-1})^2,\\
\vartheta_{10} = 2Qq^{\frac{1}{4}}\prod_{1,\infty}^{\nu}(1 + q^{2\nu})^2,\\
\vartheta_{11}' = 2\pi Qq^{\frac{1}{4}}\prod_{1,\infty}^{\nu}(1 - q^{2\nu})^2.
\end{gathered}}
Accordingly, it is possible by means of (5), § 23:
\[
\vartheta_{11}' = \pi\vartheta_{00}\vartheta_{01}\vartheta_{10}
\]
%CONT>>
